\documentclass{situationist-preprint}

\title{The Present Is Not Special}
\msid{TSM-15-a}
\msfamily{TSM-15}
\msdiscipline{Time, Space and Modality}
\mstype{Principal manuscript}
\mskeywords{present; presentness; A-theory; B-theory; indexicals; token-reflexivity}

\begin{document}
\maketitle

\begin{abstract}
Is the present moment special? A-theorists hold that one time is objectively present and that its presentness is a fundamental feature of reality. We argue that the present has no metaphysical privilege: it is the time at which a thought or utterance occurs. Once a monadic property of presentness is excluded, the only presentness left is a relation that every time bears to the thoughts and utterances occurring at it. Every time is present in this sense, relative to its own occupants, and none is present absolutely, so presentness cannot single out a time. The absence of a frame-independent present in relativity gives independent support. We answer the arguments from relief, from experience and from the question which time is really present.
\end{abstract}

\section{The problem}

On the A-theory one time, the present, is privileged: for the presentist only present things exist \citep{Markosian2004}, and for other A-theorists the present is where the objective passage of time is located. On the B-theory, \enquote{present} is like \enquote{here}: a time is present relative to a thought or utterance if the thought or utterance occurs at it \citep{Reichenbach1947,Smart1963,Mellor1998}. We argue for the second view, drawing on TSM-01.

\section{Framework}

\begin{definition}[Presentness]\label{d:pres}
A time $t$ is \emph{present relative to} a thought or utterance $u$ if $u$ occurs at $t$.
\end{definition}

\begin{definition}[Privilege]\label{d:priv}
A time is \emph{metaphysically privileged} if it has a feature that no other time has and that it has not merely relative to something located at it.
\end{definition}

\begin{assumption}[No monadic presentness]\label{a:nomon}
There is no monadic property of being present that one time has and the others lack.
\end{assumption}

Assumption~\ref{a:nomon} is the conclusion of TSM-01: a monadic presentness would make the A-series a feature of reality, and McTaggart's argument shows the A-series so understood to be incoherent.

\section{The argument}

\begin{theorem}\label{t:main}
No time is metaphysically privileged by presentness. The present is the time at which a thought or utterance occurs.
\end{theorem}

\begin{proof}
By Assumption~\ref{a:nomon}, the only presentness is that of Definition~\ref{d:pres}, which a time has relative to something located at it and which therefore does not privilege (Definition~\ref{d:priv}). Nor is it unique: every time at which a thought occurs is present relative to it, and if thoughts occur at distinct times $t_1$ and $t_2$, no time is present relative to both.
\end{proof}

TSM-16 gives independent support: a privileged present would be a global, frame-independent class of events, and relativity provides none.

\section{Objections and replies}

\begin{objection}[Relief]\label{o:relief}
After a painful experience I say \enquote{thank goodness that's over}. I am glad not that the pain ends before some date, which I could have known in advance, but that it is past now \citep{Prior1959}.
\end{objection}

\begin{reply}
What enters the attitude is an indexical way of thinking, not a privileged time. Indexical beliefs explain action and emotion as their tenseless counterparts do not \citep{Perry1979}, and a token of \enquote{the pain is over} is true if and only if it occurs after the pain. My relief is apt because it occurs then; nothing more is needed to explain why it is felt then and not earlier \citep{Mellor1998}.
\end{reply}

\begin{objection}[Experience]\label{o:exp}
Experience presents one time as present: what I see, I see as happening now.
\end{objection}

\begin{reply}
Every experience presents its own time as present, as Definition~\ref{d:pres} predicts. TSM-02 argues that the impression of passage is a feature of experience rather than of time, and the same holds of the impression that one time is singled out.
\end{reply}

\begin{objection}[Which time?]\label{o:which}
On the B-theory, which time is really present? Some time must be, or nothing would be happening.
\end{objection}

\begin{reply}
Things happen at every time. The question asks which time is present relative to no standpoint, and there is no such fact (TSM-18).
\end{reply}

\section{Conclusion}

The present is the time at which a thought or utterance occurs. Every time is present in that sense relative to its own occupants and none absolutely, so the present has no metaphysical privilege.

\begin{thebibliography}{99}\small
\bibitem[Markosian(2004)]{Markosian2004} Markosian, N. (2004). A defense of presentism. In D.~W. Zimmerman (ed.), \emph{Oxford Studies in Metaphysics}, Vol.~1, 47--82. Oxford University Press.
\bibitem[Mellor(1998)]{Mellor1998} Mellor, D.~H. (1998). \emph{Real Time II}. Routledge.
\bibitem[Perry(1979)]{Perry1979} Perry, J. (1979). The problem of the essential indexical. \emph{No{\^u}s}, 13(1), 3--21.
\bibitem[Prior(1959)]{Prior1959} Prior, A.~N. (1959). Thank goodness that's over. \emph{Philosophy}, 34(128), 12--17.
\bibitem[Reichenbach(1947)]{Reichenbach1947} Reichenbach, H. (1947). \emph{Elements of Symbolic Logic}. Macmillan.
\bibitem[Smart(1963)]{Smart1963} Smart, J.~J.~C. (1963). \emph{Philosophy and Scientific Realism}. Routledge \& Kegan Paul.
\end{thebibliography}

\end{document}
